\section{Separable Extensions}

\exercise{51.1}{Combining square and cube roots}
Choose $\alpha=\sqrt2/\sqrt[3]2=2^{1/6}$. This quotient is in the given field. Direct exponent calculations give $\alpha^3=\sqrt2$ and $\alpha^2=\sqrt[3]2$, which are polynomials in $\alpha$ with rational coefficients. Thus each original generator lies in $\Q(\alpha)$, while $\alpha$ lies in their field, proving equality.

\exercise{51.2}{Combining fourth and sixth roots}
Choose $\alpha=\sqrt[4]2/\sqrt[6]2=2^{1/12}$, an element of the given field. Then $\sqrt[4]2=\alpha^3$ and $\sqrt[6]2=\alpha^2$. These explicit polynomial expressions prove $\Q(\sqrt[4]2,\sqrt[6]2)=\Q(\alpha)$.

\exercise{51.3}{A primitive biquadratic generator}
Let $\alpha=\sqrt2+\sqrt3$, which belongs to the given field. Cubing and simplifying gives $\alpha^3=11\sqrt2+9\sqrt3$. Subtracting suitable multiples of $\alpha$ yields
\[\sqrt2=\frac{\alpha^3-9\alpha}{2},\qquad\sqrt3=\frac{11\alpha-\alpha^3}{2}.\]
Thus both generators are rational polynomials in $\alpha$, giving the desired equality.

\exercise{51.4}{A primitive generator for a complex sextic field}
Put $r=\sqrt[3]2$ and choose $\alpha=r+i$. From $(\alpha-i)^3=2$ and $i^2=-1$,
\[i=\frac{\alpha^3-3\alpha-2}{3\alpha^2-1}.\]
The denominator is nonzero: $\alpha$ has imaginary part 1 and positive real part, so its square has nonzero imaginary part. This already proves that $i,r\in\Q(\alpha)$. To meet the requested \emph{polynomial} form, exact reduction gives
\begin{align*}
i&=\frac{-12\alpha^5+9\alpha^4-40\alpha^3+78\alpha^2-78\alpha-91}{22},\\
r&=\frac{12\alpha^5-9\alpha^4+40\alpha^3-78\alpha^2+100\alpha+91}{22}.
\end{align*}
Here is a way to check these expressions without assuming the answer. Squaring the first displayed numerator equation for $i$ gives
\[m(\alpha)=0,\quad m(x)=x^6+3x^4-4x^3+3x^2+12x+5.\]
Polynomial division verifies that $(3x^2-1)$ times the displayed polynomial for $i$ is congruent to $x^3-3x-2$ modulo $m(x)$. Evaluation at $\alpha$ proves the formula; subtracting it from $\alpha$ gives the formula for $r$. Hence the field is $\Q(\alpha)$.

\exercise{51.5}{Where multiplicity is measured}
The factorization must be in $\overline F[x]$, not necessarily in $F[x]$, since $\alpha$ may lie outside $F$. For a \emph{nonzero} polynomial $f$, multiplicity $v\ge1$ means $(x-\alpha)^v\mid f$ but $(x-\alpha)^{v+1}\nmid f$ in $\overline F[x]$. The zero polynomial has no finite root multiplicity under this definition.

\exercise{51.6}{The definition of separable element}
The proposed definition is correct: an algebraic element is separable over $F$ when it has multiplicity 1 as a root of its minimal polynomial over $F$. All roots of an irreducible polynomial have the same multiplicity, so this also says the minimal polynomial has distinct roots in an algebraic closure.

\exercise{51.7}{A repeated polynomial over a perfect field}
Take $f(x)=(x^2+1)^2$. It has no rational root, and its complex roots $i,-i$ each have multiplicity 2. This does not contradict perfection of $\Q$, which concerns \emph{irreducible} polynomials. The minimal polynomial $x^2+1$ has distinct roots; the repetitions were introduced by squaring it.

\exercise{51.8}{Ten statements about separability}
\begin{enumerate}[label=(\alph*)]
\item False. Over $\Z_p(t)$, $x^p-t$ is irreducible and has just one distinct root, so adjoining that root gives an inseparable finite extension.
\item True. Every finite field is perfect.
\item True. An irreducible polynomial in characteristic zero cannot share a factor with its smaller-degree nonzero derivative.
\item False. $(x-1)^2$ has a repeated root over any field.
\item False for the same example, even over $\Q$.
\item True. This is the defining polynomial characterization of a perfect field.
\item True. Every irreducible polynomial over an algebraically closed field is linear.
\item True. An algebraic closure is algebraic over $F$ and is perfect.
\item True. Separability makes the embedding index equal the degree, and splitting makes every embedding an automorphism.
\item True. The index of a finite extension divides its degree (the quotient is its inseparable degree), and in a splitting field the index equals the group order.
\end{enumerate}

\exercise{51.9}{Separable elements are closed under field operations}
The minimal polynomial of $\beta$ over $F(\alpha)$ divides its separable minimal polynomial over $F$, so it is separable. Therefore $F(\alpha,\beta)/F(\alpha)$ and $F(\alpha)/F$ are finite separable extensions; the tower theorem makes $F(\alpha,\beta)/F$ separable. Every element of this field is then separable over $F$, including $\alpha+\beta$, $\alpha\beta$, and $\alpha/\beta$ when $\beta\ne0$.

\exercise{51.10}{A basis for a Frobenius extension}
Let $K=\Z_p(y^p)$. The relation $y^p\in K$ makes $K[y]$ spanned by $1,y,\ldots,y^{p-1}$. Being a finite-dimensional domain over $K$, it is a field: multiplication by any nonzero element is an injective linear map, hence surjective, giving an inverse. Thus $K[y]=\Z_p(y)$.
For independence, a relation $\sum_{j=0}^{p-1}r_j(y^p)y^j=0$ can be multiplied by a common denominator to give $\sum_jP_j(y^p)y^j=0$. Exponents coming from different $j$ lie in different residue classes modulo $p$, so no cancellation between the groups is possible; all $P_j$ vanish. Thus the displayed spanning set is a basis, and the degree is $p$. Setting $t=y^p$, the polynomial $x^p-t$ has $y$ as a root and the same degree as its minimal polynomial. Hence it is irreducible over $\Z_p(t)$.

\exercise{51.11}{Algebraic extensions of a perfect field}
Let $\alpha$ be algebraic over $E$. Since $E/F$ is algebraic, $\alpha$ is algebraic over $F$ as well. The minimal polynomial of $\alpha$ over the perfect field $F$ is separable. Its minimal polynomial over $E$ divides it in $E[x]$, so also has no repeated roots. Thus every element algebraic over $E$ is separable over $E$, proving $E$ perfect.

\exercise{51.12}{Transitivity for infinite separable extensions}
Take $\alpha\in K$. Its separable minimal polynomial over $E$ uses finitely many coefficients $c_1,\ldots,c_m\in E$. Set $E_0=F(c_1,\ldots,c_m)$; this is a finite separable extension of $F$. The same polynomial is irreducible over $E_0$, because a factorization there would be one over $E$, and it remains separable. Thus $E_0(\alpha)/E_0$ is finite separable. The finite tower theorem gives $E_0(\alpha)/F$ separable, so $\alpha$ is separable over $F$. This holds for every $\alpha\in K$, proving transitivity without assuming the entire tower finite.

\exercise{51.13}{The separable closure inside an algebraic extension}
Let $S$ be the elements of $E$ separable over $F$. It contains $F$, including 0 and 1. Exercise 51.9 shows it is closed under sums, products, and quotients by nonzero elements; additive inverses are included by multiplying by $-1\in F$. Hence $S$ is a subfield. It consists exactly of the separable elements, so is the largest intermediate extension separable over $F$, called the separable closure of $F$ in $E$.

\exercise{51.14}{Frobenius generates a finite-field Galois group}
\begin{enumerate}[label=(\alph*)]
\item In a field of $p^n$ elements, $a^{p^n}=a$ for every $a$, so $\sigma_p^n=1$. If $\sigma_p^d=1$ for $0<d<n$, all $p^n$ elements would be roots of $x^{p^d}-x$, whose degree $p^d$ is smaller than $p^n$, impossible. Thus its order is $n$.
\item The extension has degree $n$, is separable because finite fields are perfect, and is a splitting extension since it splits $x^{p^n}-x$. Its Galois group therefore has order $n$. The $n$ distinct powers of Frobenius already account for that many elements, so the full group is cyclic with this generator.
\end{enumerate}

\exercise{51.15}{The additive derivative homomorphism}
\begin{enumerate}[label=(\alph*)]
\item If $f=\sum a_jx^j$ and $g=\sum b_jx^j$, differentiating coefficient by coefficient gives $(f+g)'=\sum j(a_j+b_j)x^{j-1}=f'+g'$. This proves it is a homomorphism of additive groups; it is generally not a ring homomorphism.
\item In characteristic zero, $ja_j=0$ for $j>0$ implies $a_j=0$. Thus the kernel is the constant polynomials $F$.
\item In characteristic $p$, $j$ vanishes as a field scalar precisely when $p\mid j$. Hence the kernel is $F[x^p]$, the polynomials all of whose nonzero exponents are multiples of $p$. Their coefficients need not themselves be $p$th powers.
\end{enumerate}

\exercise{51.16}{Formal derivative rules}
\begin{enumerate}[label=(\alph*)]
\item For $c\in F$, coefficientwise differentiation gives $(cf)'=\sum_jjca_jx^{j-1}=cf'$.
\item First check $(x^ix^j)'=(i+j)x^{i+j-1}=x^i(jx^{j-1})+(ix^{i-1})x^j$. Expand arbitrary polynomials into monomials and use additivity and (a) to obtain $(fg)'=fg'+f'g$.
\item Induct on $m$. For $m=1$ the assertion is immediate. If $(f^m)'=mf^{m-1}f'$, the product rule gives $(f^{m+1})'=f^mf'+(f^m)'f=(m+1)f^mf'$. The integers are interpreted as elements of $F$, so this proof also works in positive characteristic.
\end{enumerate}

\exercise{51.17}{Multiplicity and vanishing derivative}
In $\overline F[x]$, write $f=(x-\alpha)^vg$ with $g(\alpha)\ne0$. If $v=1$, the product rule gives $f'(\alpha)=g(\alpha)\ne0$. If $v>1$, both terms of
\[f'=v(x-\alpha)^{v-1}g+(x-\alpha)^vg'\]
vanish at $\alpha$, including when the scalar $v$ is zero in the field. Thus $f'(\alpha)=0$ exactly when the root has multiplicity greater than 1.

\exercise{51.18}{Characteristic zero forces irreducible separability}
Let $p\in F[x]$ be nonconstant and irreducible in characteristic zero. Its derivative is nonzero and has smaller degree. If a root $\alpha$ were repeated, it would be a root of $p'$ too, so its minimal polynomial $p$ would divide $p'$. The degree inequality rules this out. Therefore all roots of $p$ are simple, proving separability.

\exercise{51.19}{The derivative test for an inseparable irreducible}
For an irreducible $q$, a repeated root implies $q\mid q'$ by minimality. Since $\deg q'<\deg q$ unless $q'=0$, this is equivalent to $q'=0$. In characteristic $p$, coefficientwise differentiation gives $q'=0$ exactly when each exponent with nonzero coefficient is divisible by $p$. Conversely a nonconstant polynomial with zero derivative cannot have a simple root by Exercise 51.17. Thus the criterion holds in both directions.

\exercise{51.20}{Common factors over the algebraic closure}
By Exercise 51.17 a repeated root is exactly a common root of $f$ and $f'$. Over the algebraically closed field $\overline F$, a common root $\alpha$ supplies the common linear factor $x-\alpha$. Conversely a nonconstant common factor has a root in the closure, which is a common root of $f,f'$, hence repeated in $f$. This proves the requested equivalence for nonzero $f$.

\exercise{51.21}{Bringing the gcd test back to the coefficient field}
If $f,f'$ have no nonconstant common divisor in $F[x]$, their gcd is 1, so Bezout's identity gives $Af+Bf'=1$ with $A,B\in F[x]$. This same identity holds in $\overline F[x]$, excluding any nonconstant common factor there. Conversely a nonconstant common factor in $F[x]$ remains one over the closure. Combining with Exercise 51.20 proves the equivalence over $F$ itself. The argument includes nonzero constant $f$, whose gcd with its zero derivative is a unit.

\exercise{51.22}{A root-free computational test}
For a nonzero polynomial, form its formal derivative and run the Euclidean algorithm in $F[x]$ on $f,f'$. Normalize the final nonzero remainder to a monic gcd $d$. A repeated root exists exactly when $\deg d>0$. In particular, if $f$ is nonconstant and $f'=0$, the gcd is the monic associate of $f$, so there is repetition. If $f$ is a nonzero constant, the gcd has degree 0 and there are no roots. No numerical root approximation is required.
